                         \qquad g\in\mathcal P_1.          \tag{ND.9}
\] Indeed, put \(x=|z_0|^2\), \(y=|z_2|^2\), \(w=|z_1|^2+|z_3|^2\). Equation (ND.8) and the reverse triangle inequality give \(\Gamma_T(g)\geq8\pi^2[(x-2y)^2+4y^2+w^2]\). Cauchy--Schwarz gives \((x+y+w)^2\leq(17/4)[(x-2y)^2+4y^2+w^2]\), and (ND.4) proves (ND.9). The sharper even-symbol constant from the earlier project is not silently transferred to arbitrary Fourier phases.

\subsubsection{A constructive margin for the radius-eight certificate}\label{nondegeneracy:a-constructive-margin-for-the-radius-eight-certificate}

Compactness already gives some \(c_8>0\) with \(\Gamma_T(g)\geq c_8\|g\|_2^4\). There is also a completely explicit, deliberately very poor lower bound: \[
 \boxed{\displaystyle
 \Gamma_T(g)\geq \frac{\pi^2}{41472}\,2^{-2^{715}}\|g\|_2^4,
                         \qquad g\in\mathcal P_8.}          \tag{ND.10}
\] This number has no useful interpretation as a physical crossover estimate; its purpose is to remove an unspecified compactness constant. It is not claimed to approximate the sharp minimum.

Here is the quantitative compilation rule. Normalize \(\max_i|z_i|=1\) and suppose \(|Q_j|\leq\delta\). Replace an exact zero fact by an upper bound \(2^L\delta^{p-kb}\) for its absolute value. The symbol \(b\) appears only under a temporary assumption \(|z_i|\geq\delta^b\). Every deleted monomial is bounded by a known coefficient bound times \(|z_j|\leq1\), or by a known product bound. In the residual row, use the triangle inequality, the minimum of the base exponents, the maximum of their \(b\) slopes, and a power-of-two upper bound on the sum of absolute integer coefficients. P records the resulting product bound. S takes a square root, rounding its power-of-two coefficient upward. D subtracts \(b\) from the exponent using the explicit nonzero lower bound.

At the end of a refutation child, the pivot has bound \(2^L\delta^{p-kb}\) with \(p>0\). Choose \(b=p/[2(k+1)]\) and require \(\delta\leq2^{-B}\) with \(B\geq2(L+1)/p\). Then this bound is strictly less than \(\delta^b\), contradicting the temporary assumption. Therefore the pivot has the unconditional bound \(|z_i|<\delta^b\). Proceed through the root proof and finally enlarge \(B\) so that every coefficient is at most \(1/2\). This contradicts the normalization. All exponents and comparisons in the compiler are exact rational arithmetic; there is no numerical interval test on the coefficients of \(g\).

For the maintained radius-eight proof the generated integer \(B\) has 714 binary digits, so the simple upper bound \(B\leq2^{714}\) suffices. The preceding argument gives \(\max_j|Q_j|>2^{-B}\) on \(\max_i|z_i|=1\). Since \(\|g\|_2^4\leq4n_R^2(\max_i|z_i|)^4\), (ND.6) yields \(\Gamma_T(g)\geq[\pi^2/(2n_R^2)]2^{-2B}\|g\|_2^4\); with \(n_8=144\) this implies (ND.10). The exact, much less conveniently printed integer and every intermediate rational bound are retained in the margin report. The original margin report was not deposited; it has been reconstructed from the deposited certificate by an exact-rational-arithmetic implementation of the compilation rule above (\path{computations/verification/resonance_nondegeneracy_margin.py}), and the reconstructed report is deposited at \path{evidence/executions/20261006/corpus-revision/final/checks/certificates-mpuvq_28/resonance_nondegeneracy_margin.json}. The reconstructed integer has 712 binary digits, so \(B\leq2^{714}\) is confirmed and the explicit constant in (ND.10) is recheckable from deposited materials; two independent implementations of the compiler produce the identical reconstructed integer. The unconditional compactness bound \(\Gamma_T(g)\geq c_8\|g\|_2^4\) for some \(c_8>0\) does not depend on the margin report.

\subsection{An invariant Hessian obstruction for every hyperbolic matrix}\label{nondegeneracy:an-invariant-hessian-obstruction-for-every-hyperbolic-matrix}

The finite certificate settles a large but fixed frequency box. A separate analytic criterion has no bandwidth restriction and is not specific to T. Write \[
 A=\begin{pmatrix}a&b\\c&d\end{pmatrix},\quad t=a+d,\quad
 \Delta=t^2-4>0,\quad
 B_A=\begin{pmatrix}c&(d-a)/2\\(d-a)/2&-b\end{pmatrix}.
\] The integral quadratic form \[
 q_A(v)=v^{\mathsf T}B_Av=\det(v,Av)
\] is A-invariant. Define, for a symmetric matrix H, \[
